% ============================================================
%  S&P 500 Stock Price Forecasting using Deep Learning
%  COMS 4030A / COMS 7047A -- Adaptive Computation and Machine Learning
%  University of Cape Town -- May 2026
% ============================================================
\documentclass[12pt, a4paper]{article}

% -- Encoding & fonts ----------------------------------------
\usepackage[T1]{fontenc}
\usepackage[utf8]{inputenc}
\usepackage{mathptmx}
\usepackage[scaled=0.92]{helvet}
\usepackage{microtype}

% -- Page geometry -------------------------------------------
\usepackage[
  a4paper,
  top=2.5cm, bottom=2.5cm,
  left=3cm,  right=2.5cm,
  headheight=24pt
]{geometry}

% -- Colour --------------------------------------------------
\usepackage[dvipsnames, table]{xcolor}
\definecolor{uctnavy}{HTML}{003366}
\definecolor{uctgold}{HTML}{C9902A}
\definecolor{tabhead}{HTML}{003366}
\definecolor{tabrow}{HTML}{EEF2F8}
\definecolor{tabrowalt}{HTML}{F8F9FC}

% -- Headers & footers ---------------------------------------
\usepackage{fancyhdr}
\pagestyle{fancy}
\fancyhf{}
\renewcommand{\headrulewidth}{0.5pt}
\renewcommand{\headrule}{\hbox to\headwidth{\color{uctnavy}\leaders\hrule height \headrulewidth\hfill}}
\renewcommand{\footrulewidth}{0.5pt}
\renewcommand{\footrule}{\hbox to\headwidth{\color{uctnavy}\leaders\hrule height \footrulewidth\hfill}}
\fancyhead[L]{\small\color{uctnavy}\textit{COMS 4030A / 7047A --- Adaptive Computation and Machine Learning}}
\fancyhead[R]{\small\color{uctnavy}\textit{University of Cape Town}}
\fancyfoot[L]{\small\color{gray}S\&P 500 Stock Price Forecasting $|$ May 2026}
\fancyfoot[C]{\small\color{uctnavy}\textbf{\thepage}}
\fancyfoot[R]{\small\color{gray}Ebrahim \& Tayob $|$ UCT}

\fancypagestyle{plain}{%
  \fancyhf{}
  \renewcommand{\headrulewidth}{0pt}
  \renewcommand{\footrulewidth}{0pt}
}

% -- Section formatting --------------------------------------
\usepackage{titlesec}
\titleformat{\section}
  {\color{uctnavy}\large\bfseries\raggedright}
  {\thesection.}{0.6em}{}
  [\vspace{-4pt}\color{uctnavy}\rule{\linewidth}{0.5pt}\vspace{2pt}]

\titleformat{\subsection}
  {\color{uctnavy}\normalsize\bfseries\raggedright}
  {\thesubsection}{0.6em}{}

\titlespacing*{\section}{0pt}{18pt}{8pt}
\titlespacing*{\subsection}{0pt}{12pt}{4pt}

% -- Tables --------------------------------------------------
\usepackage{booktabs}
\usepackage{tabularx}
\usepackage{array}
\newcolumntype{L}[1]{>{\raggedright\arraybackslash}p{#1}}
\newcolumntype{C}[1]{>{\centering\arraybackslash}p{#1}}
\newcolumntype{X}{>{\raggedright\arraybackslash}X}

\usepackage{caption}
\captionsetup[table]{
  labelfont={bf, color=uctnavy},
  textfont=it,
  position=above,
  skip=4pt
}
\captionsetup[figure]{
  labelfont={bf, color=uctnavy},
  textfont=it,
  position=below,
  skip=6pt
}

% -- Graphics ------------------------------------------------
\usepackage{graphicx}
\graphicspath{{figures/}}

% -- Lists ---------------------------------------------------
\usepackage{enumitem}
\setlist[itemize]{leftmargin=1.5em, itemsep=2pt, parsep=0pt, topsep=4pt}

% -- Hyperlinks ----------------------------------------------
\usepackage[
  colorlinks=true,
  linkcolor=uctnavy,
  urlcolor=uctnavy,
  citecolor=uctnavy,
  bookmarks=true
]{hyperref}

% -- Paragraph spacing ---------------------------------------
\setlength{\parindent}{0pt}
\setlength{\parskip}{6pt}
\linespread{1.15}

% -- Utilities -----------------------------------------------
\usepackage{amsmath}
\usepackage{float}

% ============================================================
\begin{document}
\thispagestyle{plain}

% -- TITLE PAGE ----------------------------------------------
\begin{titlepage}
  \centering
  \vspace*{3cm}

  {\color{uctnavy}\rule{\linewidth}{1.5pt}}
  \vspace{0.6cm}

  {\color{uctnavy}\LARGE\bfseries
    S\&P 500 Stock Price Forecasting\\[6pt]
    using Deep Learning\par}

  \vspace{0.5cm}
  {\color{uctnavy}\rule{\linewidth}{0.5pt}}
  \vspace{1.2cm}

  {\large\itshape
    COMS 4030A / COMS 7047A\\[4pt]
    Adaptive Computation and Machine Learning\par}

  \vspace{0.5cm}
  {\large University of Cape Town\par}

  \vspace{2.5cm}

  {\normalsize\bfseries\color{uctnavy} Team Members\par}
  \vspace{0.4cm}
  {\large
    Zeeshaan Ebrahim \quad {\normalsize\color{gray}(2543004)}\\[6pt]
    Tazeem Tayob \quad {\normalsize\color{gray}(2659956)}\par}

  \vfill
  {\large May 2026\par}
  \vspace{1cm}
  {\color{uctnavy}\rule{\linewidth}{1.5pt}}
\end{titlepage}

% -- TABLE OF CONTENTS ----------------------------------------
\newpage
\tableofcontents
\newpage

% ============================================================
\section{Dataset Description \& Objectives}

\textbf{Source:} Kaggle --- S\&P 500 Top 10 Stocks OHLCV (2010--2026). The dataset contains split-adjusted daily trading data for nine S\&P 500 technology companies, spanning over 15 years of market history.

\vspace{8pt}
\begin{table}[H]
\caption{Dataset Columns}
\label{tab:dataset-cols}
\rowcolors{2}{tabrow}{tabrowalt}
\begin{tabularx}{\textwidth}{L{3.2cm} X}
\rowcolor{tabhead}
\color{white}\textbf{Column} & \color{white}\textbf{Description} \\
\midrule
Date        & Trading date (YYYY-MM-DD) \\
Ticker      & Stock symbol (9 unique tickers) \\
Open        & Opening price (split-adjusted) \\
High        & Daily high price \\
Low         & Daily low price \\
Close       & Closing price \\
Adj\_Close  & Adjusted closing price (primary target) \\
Volume      & Number of shares traded \\
\bottomrule
\end{tabularx}
\end{table}

\textbf{Tickers:} AAPL, AMZN, AVGO, GOOG, GOOGL, META, MSFT, NVDA, TSLA

\textbf{Dataset statistics:} 35,765 rows $|$ 9 stocks $|$ Date range: 2010-01-04 to 2026-02-13 $|$ Missing values: 0

\subsection{Objectives}

\begin{itemize}
  \item Build a single unified deep learning model to forecast the next trading day's adjusted closing price for any of the nine stocks simultaneously.
  \item Compare four architectures: Stacked LSTM, Bidirectional LSTM with Attention, CNN-LSTM Hybrid, and custom Transformer Encoder.
  \item Investigate how window size, feature richness, dropout, prediction horizon, and L2 regularisation affect forecast accuracy.
\end{itemize}

% ============================================================
\section{Preprocessing \& Data Splits}

\subsection{Feature Engineering}

17 additional technical indicators were computed per stock using only information available up to time $t$ (no lookahead bias). All rolling calculations were performed per-ticker to prevent cross-stock contamination.

\vspace{8pt}
\begin{table}[H]
\caption{Engineered Features}
\label{tab:features}
\rowcolors{2}{tabrow}{tabrowalt}
\begin{tabularx}{\textwidth}{L{3.8cm} X}
\rowcolor{tabhead}
\color{white}\textbf{Feature Group} & \color{white}\textbf{Features} \\
\midrule
Trend (SMA)   & SMA\_5, SMA\_20, SMA\_50 \\
Trend (EMA)   & EMA\_12, EMA\_26, MACD \\
Momentum      & RSI\_14 \\
Volatility    & BB\_upper, BB\_lower, BB\_bandwidth \\
Volume        & Volume\_MA\_10 \\
Returns       & Daily\_Return, Volatility\_10 \\
Identity      & Ticker one-hot encoding (9 dimensions) \\
\bottomrule
\end{tabularx}
\end{table}

\textbf{Total: 28 features per timestep.}

\subsection{Normalisation}

MinMaxScaler was fitted exclusively on the training slice of each stock to prevent data leakage. The same fitted scaler was then applied to the validation and test slices. Scalers were stored per-ticker to ensure correct inverse transformation during evaluation.

\subsection{Data Splits}

\vspace{4pt}
\begin{table}[H]
\caption{Data Splits}
\label{tab:splits}
\rowcolors{2}{tabrow}{tabrowalt}
\begin{tabularx}{\textwidth}{L{2.2cm} L{4.2cm} C{2.4cm} X}
\rowcolor{tabhead}
\color{white}\textbf{Split} & \color{white}\textbf{Date Range} & \color{white}\textbf{Sequences} & \color{white}\textbf{Purpose} \\
\midrule
Train      & 2010-01-04 to 2022-12-31 & 27,746 & Fit models and learn weights \\
Validation & 2023-01-01 to 2024-06-30 &  2,826 & Hyperparameter tuning and early stopping \\
Test       & 2024-07-01 to 2026-02-13 &  3,132 & Final evaluation only (never seen during training) \\
\bottomrule
\end{tabularx}
\end{table}

\subsection{Sequence Creation}

A sliding window of 60 trading days ($\approx$3 months) was used. For each sample, $X =$ features for days $[t-60,\, t-1]$ and $y =$ normalised Adj\_Close at day $t$. Sequences from all nine stocks were stacked into a single shared dataset.

% ============================================================
\section{Model Selection \& Implementation}

All models were implemented in PyTorch 2.2.2 and trained with the Adam optimiser ($\text{lr}=10^{-3}$), MSE loss, gradient clipping (max norm $= 1.0$), EarlyStopping (patience $= 10$), and ReduceLROnPlateau (factor $= 0.5$, patience $= 5$).

\subsection{Model 1 --- Stacked LSTM (Baseline)}

A two-layer stacked LSTM with dropout regularisation serves as the baseline.

\textbf{Architecture:} Input$(60, 28) \to$ LSTM$(128,\;\text{return\_seq}=\text{True}) \to$ Dropout$(0.2) \to$ LSTM$(64) \to$ Dropout$(0.2) \to$ Linear$(32) \to$ ReLU $\to$ Linear$(1)$

\textbf{Parameters:} 132,673 \quad $|$ \quad \textbf{Best validation loss:} 0.070183

\subsection{Model 2 --- Bidirectional LSTM with Bahdanau Attention}

Bidirectional LSTMs capture both forward and backward temporal context within the 60-day window. A custom Bahdanau-style attention mechanism computes a weighted sum of hidden states, allowing the model to focus on the most predictive timesteps. The final context vector is passed through a fully connected head.

\textbf{Parameters:} 330,946 \quad $|$ \quad \textbf{Best validation loss:} 0.115797

\subsection{Model 3 --- CNN-LSTM Hybrid}

Temporal causal convolutions act as local pattern extractors detecting short-term momentum shifts, while the LSTM captures long-range sequential dependencies. Causal padding (left-padding only) prevents any future information leakage into the convolutional feature maps.

\textbf{Parameters:} 51,169 \quad $|$ \quad \textbf{Best validation loss:} 0.012034

\subsection{Model 4 --- Custom Transformer Encoder}

Built entirely from scratch using PyTorch primitives following the Vaswani et al.\ (2017) encoder design with multi-head self-attention, position-wise feed-forward networks, layer normalisation, and residual connections. Sinusoidal positional encodings were applied to the input embeddings.

\textbf{Parameters:} 76,865 \quad $|$ \quad \textbf{Best validation loss:} 0.122278

% ============================================================
\section{Experiments \& Fine-Tuning}

Five controlled experiments were conducted using the LSTM as a baseline (except Experiment 5 which used the Transformer).

\subsection{Experiment 1 --- Window Size}

The sequence length was varied across 20, 30, and 60 trading days to assess the impact of lookback horizon on forecast accuracy.

\vspace{4pt}
\begin{table}[H]
\caption{Window Size vs.\ Validation RMSE}
\label{tab:window}
\rowcolors{2}{tabrow}{tabrowalt}
\begin{tabularx}{\textwidth}{C{5cm} C{5cm}}
\rowcolor{tabhead}
\color{white}\textbf{Window Size} & \color{white}\textbf{Val RMSE} \\
\midrule
20 days & 0.1599 \\
30 days & 0.1952 \\
\textbf{60 days (best)} & \textbf{0.0749} \\
\bottomrule
\end{tabularx}
\end{table}

\textbf{Result:} A 60-day window significantly outperformed shorter windows, capturing quarterly price trends and momentum cycles invisible at shorter horizons.

\subsection{Experiment 2 --- Feature Ablation}

OHLCV-only (6 features) was compared against the full technical indicator set (28 features).

\vspace{4pt}
\begin{table}[H]
\caption{Feature Set vs.\ Validation RMSE}
\label{tab:feat-ablation}
\rowcolors{2}{tabrow}{tabrowalt}
\begin{tabularx}{\textwidth}{X C{4cm}}
\rowcolor{tabhead}
\color{white}\textbf{Feature Set} & \color{white}\textbf{Val RMSE} \\
\midrule
OHLCV-only (6 features) & Higher RMSE \\
\textbf{Full features with indicators (28)} & \textbf{Lower RMSE (winner)} \\
\bottomrule
\end{tabularx}
\end{table}

\textbf{Result:} Technical indicators (MACD, RSI, Bollinger Bands) provided meaningful additional signal, reducing RMSE compared to raw price and volume data alone.

\subsection{Experiment 3 --- Dropout Rate}

Three dropout rates were tested on the LSTM.

\vspace{4pt}
\begin{table}[H]
\caption{Dropout Rate vs.\ Validation RMSE}
\label{tab:dropout}
\rowcolors{2}{tabrow}{tabrowalt}
\begin{tabularx}{\textwidth}{C{5cm} C{5cm}}
\rowcolor{tabhead}
\color{white}\textbf{Dropout Rate} & \color{white}\textbf{Val RMSE} \\
\midrule
\textbf{0.1 (best)} & \textbf{0.1325} \\
0.2                 & 0.2322 \\
0.3                 & 0.3171 \\
\bottomrule
\end{tabularx}
\end{table}

\textbf{Result:} A dropout rate of 0.1 achieved the lowest validation RMSE. Higher dropout rates caused underfitting.

\subsection{Experiment 4 --- Prediction Horizon}

1-day and 5-day ahead forecasting were compared on the LSTM.

\vspace{4pt}
\begin{table}[H]
\caption{Prediction Horizon vs.\ Validation RMSE}
\label{tab:horizon}
\rowcolors{2}{tabrow}{tabrowalt}
\begin{tabularx}{\textwidth}{C{5cm} C{5cm}}
\rowcolor{tabhead}
\color{white}\textbf{Horizon} & \color{white}\textbf{Val RMSE} \\
\midrule
\textbf{1 day ahead (winner)} & \textbf{Lower RMSE} \\
5 days ahead & Higher RMSE ($\approx$doubled) \\
\bottomrule
\end{tabularx}
\end{table}

\textbf{Result:} The 5-day horizon was significantly harder, with RMSE roughly doubling, indicating rapid error accumulation over multiple timesteps.

\subsection{Experiment 5 --- L2 Regularisation (Transformer)}

Weight decay was applied to the Transformer optimiser across three settings.

\vspace{4pt}
\begin{table}[H]
\caption{Weight Decay vs.\ Validation RMSE (Transformer)}
\label{tab:l2}
\rowcolors{2}{tabrow}{tabrowalt}
\begin{tabularx}{\textwidth}{C{5cm} C{5cm}}
\rowcolor{tabhead}
\color{white}\textbf{Weight Decay} & \color{white}\textbf{Val RMSE} \\
\midrule
0 (none)                    & 0.3389 \\
\textbf{$10^{-4}$ (best)}  & \textbf{0.3250} \\
$10^{-3}$                   & 0.3484 \\
\bottomrule
\end{tabularx}
\end{table}

\textbf{Result:} Mild L2 regularisation ($10^{-4}$) slightly improved generalisation. Strong regularisation ($10^{-3}$) caused underfitting.

% ============================================================
\section{Training Visualisations}

\begin{figure}[H]
  \centering
  \includegraphics[width=\textwidth]{fig1}
  \caption{Normalised adjusted closing prices for all 9 stocks (2010--2026).}
  \label{fig:norm-prices}
\end{figure}

\begin{figure}[H]
  \centering
  \includegraphics[width=\textwidth]{fig2}
  \caption{Cross-stock return correlation heatmap.}
  \label{fig:corr}
\end{figure}

\begin{figure}[H]
  \centering
  \includegraphics[width=\textwidth]{fig3}
  \caption{Daily trading volume distribution by stock.}
  \label{fig:volume}
\end{figure}

\begin{figure}[H]
  \centering
  \includegraphics[width=\textwidth]{fig4}
  \caption{Training and validation MSE loss curves for all four models.}
  \label{fig:loss}
\end{figure}

\begin{figure}[H]
  \centering
  \includegraphics[width=\textwidth]{fig5}
  \caption{Experiment results: window size, feature ablation, and dropout rate.}
  \label{fig:exp135}
\end{figure}

\begin{figure}[H]
  \centering
  \includegraphics[width=\textwidth]{fig6}
  \caption{Experiment results: prediction horizon and L2 regularisation.}
  \label{fig:exp45}
\end{figure}

\begin{figure}[H]
  \centering
  \includegraphics[width=\textwidth]{fig7}
  \caption{BiLSTM attention weight heatmap over the 60-day lookback window.}
  \label{fig:attention}
\end{figure}

% ============================================================
\section{Test Results \& Analysis}

All predictions were inverse-transformed from normalised $[0, 1]$ space back to actual dollar prices before computing metrics.

\vspace{4pt}
\begin{table}[H]
\caption{Test Set Performance Across All Models}
\label{tab:results}
\rowcolors{2}{tabrow}{tabrowalt}
\begin{tabularx}{\textwidth}{L{4.2cm} C{2.2cm} C{2.2cm} C{2.4cm} C{2.8cm}}
\rowcolor{tabhead}
\color{white}\textbf{Model} &
\color{white}\textbf{RMSE (\$)} &
\color{white}\textbf{MAE (\$)} &
\color{white}\textbf{MAPE (\%)} &
\color{white}\textbf{Dir.\ Acc.\ (\%)} \\
\midrule
LSTM                 & 79.953  & 55.399 & 19.857 & 49.38 \\
BiLSTM + Attention   & 98.377  & 69.123 & 23.934 & 50.40 \\
\textbf{CNN-LSTM}    & \textbf{41.827}  & \textbf{26.785} & \textbf{10.278} & \textbf{51.17} \\
Transformer          & 101.240 & 71.212 & 24.502 & 50.46 \\
\bottomrule
\end{tabularx}
\end{table}

\textbf{Best model:} CNN-LSTM with RMSE \$41.83, MAE \$26.79, MAPE 10.28\%, Directional Accuracy 51.17\%.

\begin{figure}[H]
  \centering
  \includegraphics[width=\textwidth]{fig8}
  \caption{Test set performance comparison --- RMSE, MAE, MAPE across all models.}
  \label{fig:perf}
\end{figure}

\begin{figure}[H]
  \centering
  \includegraphics[width=\textwidth]{fig9}
  \caption{Actual vs.\ predicted Adj\_Close for AAPL, NVDA, TSLA (test set).}
  \label{fig:actual-pred}
\end{figure}

\begin{figure}[H]
  \centering
  \includegraphics[width=\textwidth]{fig10}
  \caption{Residual error distribution per stock --- CNN-LSTM (best model).}
  \label{fig:residuals}
\end{figure}

\subsection{Analysis}

The CNN-LSTM hybrid achieved the best performance across all metrics by a significant margin. Its RMSE of \$41.83 was roughly half that of the LSTM (\$79.95) and less than half that of the Transformer (\$101.24). The local pattern extraction performed by the causal convolutional layers proved highly effective at capturing short-term momentum shifts before they are passed to the LSTM for longer-range modelling.

The Bidirectional LSTM with Attention underperformed relative to expectations. While attention mechanisms are theoretically beneficial, the added model complexity (330K parameters, $2.5\times$ more than the LSTM baseline) led to overfitting despite dropout regularisation.

The Transformer Encoder was the weakest performer (\$101.24 RMSE). Without pre-training on large datasets, Transformers tend to require substantially more data to generalise---the approximately 27,000 training sequences available were insufficient to exploit the architecture's full capacity.

All four models achieved directional accuracy close to 50\%, slightly better than random chance. This is consistent with the Efficient Market Hypothesis, which posits that short-term price direction is inherently difficult to predict from historical price data alone.

Residual distributions for all stocks are approximately centred around zero, indicating no systematic bias. However, residuals are wider for high-volatility stocks (NVDA, TSLA) compared to more stable equities (AAPL, MSFT), reflecting the greater difficulty of forecasting highly volatile assets.

% ============================================================
\section{Conclusions}

Four deep learning architectures were trained on a unified dataset of nine S\&P 500 technology stocks to forecast next-day adjusted closing prices. The CNN-LSTM hybrid was the clear best performer, achieving the lowest RMSE, MAE, and MAPE across all nine stocks.

\subsection{What Worked}

\begin{itemize}
  \item Engineering technical indicators (MACD, RSI, Bollinger Bands) reduced RMSE compared to raw OHLCV data alone.
  \item A 60-day lookback window significantly outperformed 20- and 30-day windows, capturing quarterly trend cycles.
  \item Causal convolutional feature extraction (CNN-LSTM) was more effective than pure recurrent or attention-based approaches.
  \item ReduceLROnPlateau scheduling and early stopping prevented overfitting and reduced training time.
\end{itemize}

\subsection{What Did Not Work}

\begin{itemize}
  \item The 5-day prediction horizon was significantly harder, with RMSE roughly doubling.
  \item High dropout rates (0.3) degraded performance, suggesting underfitting beyond this threshold.
  \item Strong L2 regularisation (weight\_decay $= 10^{-3}$) caused underfitting in the Transformer.
  \item The Bidirectional LSTM with Attention overfitted despite dropout regularisation.
\end{itemize}

\subsection{Limitations}

\begin{itemize}
  \item No fundamental data (earnings, macroeconomic indicators) were used; models rely solely on technical price history.
  \item Financial time-series are non-stationary; concept drift events (e.g.\ market crashes) pose challenges not addressed here.
  \item Directional accuracy near 50\% limits practical trading utility.
\end{itemize}

\subsection{Future Work}

\begin{itemize}
  \item \textbf{Cross-stock attention:} allow each stock to attend to signals from the other eight.
  \item \textbf{Multi-step forecasting:} recursive forecasting with uncertainty quantification via Monte Carlo Dropout.
  \item \textbf{Ensemble methods:} combine all four architectures for lower-variance predictions.
  \item \textbf{Richer features:} incorporate fundamental and macroeconomic indicators.
\end{itemize}

\end{document}
